\documentclass[12pt]{article}
\usepackage[T1]{fontenc}
\usepackage[utf8]{inputenc}
\usepackage{lmodern}
\usepackage{textcomp}
\usepackage{enumitem}
\usepackage{geometry}
\geometry{margin=1in}
\begin{document}
\begin{center}
Answer Key - Geography - Graduate Level Assessment 13 \
\small (Answers are ordered exactly as in the question paper.)
\end{center}

\section*{Multiple Choice}
\begin{enumerate}[label=\arabic*.]
\item \textbf{Answer:} B
\\ \textit{Options:} A) 9\%; B) 19\%; C) 39\%; D) 59\%
\\ \textit{Note:} The correct answer of 19\% reflects the findings from a 2015 survey that specifically assessed the prevalence of physical fights among school-aged children in China, highlighting a significant issue regarding youth violence in educational settings.
\item \textbf{Answer:} B
\\ \textit{Options:} A) strongly supported by most studies; B) supported mainly by cross-section, not time-series studies; C) supported mainly by time-series, not cross-section studies; D) generally repudiated by empirical studies
\\ \textit{Note:} The inverted U hypothesis, which posits that inequality initially rises and subsequently falls with economic development, is primarily substantiated by cross-sectional studies that capture data at a single point in time across different regions, rather than time-series studies that track changes over time.
\item \textbf{Answer:} B
\\ \textit{Options:} A) 0.50\%; B) 2\%; C) 6\%; D) 12\%
\\ \textit{Note:} The correct answer is based on China's reported military expenditure as a percentage of its GDP, which was approximately 2\% in 2017, reflecting the country's strategic investment in defense amid its rapid economic growth.
\item \textbf{Answer:} B
\\ \textit{Options:} A) The global growth rate was four times as high 50 years ago as it is in 2020.; B) The global growth rate was two times as high 50 years ago as it is in 2020.; C) The global growth rate is two times as high as it is in 2020.; D) The global growth rate is four times as high today as it is in 2020.
\\ \textit{Note:} The correct answer is based on historical demographic data indicating that global population growth rates peaked in the late 20 th century, specifically around the 1970 s, and have since experienced a significant decline by 2020.
\item \textbf{Answer:} C
\\ \textit{Options:} A) 38\%; B) 58\%; C) 78\%; D) 98\%
\\ \textit{Note:} The correct answer of 78\% reflects the significant prevalence of physical punishment in Indian schools as documented in studies conducted around 2009, highlighting a concerning aspect of educational practices and child welfare in the country.
\end{enumerate}

\section*{True / False}
\begin{enumerate}[label=\arabic*.]
\setcounter{enumi}{5}
\item \textbf{Answer:} False
\\ \textit{Options:} A) True; B) False
\\ \textit{Note:} The correct answer is: National Central Museum
\item \textbf{Answer:} True
\\ \textit{Options:} A) True; B) False
\\ \textit{Note:} According to the passage: handicrafts
\item \textbf{Answer:} False
\\ \textit{Options:} A) True; B) False
\\ \textit{Note:} The correct answer is: adult contemporary
\item \textbf{Answer:} True
\\ \textit{Options:} A) True; B) False
\\ \textit{Note:} According to the passage: 1526
\item \textbf{Answer:} True
\\ \textit{Options:} A) True; B) False
\\ \textit{Note:} According to the passage: Last Glacial Maximum
\end{enumerate}

\section*{Long Form Response}
\begin{enumerate}[label=\arabic*.]
\setcounter{enumi}{10}
\item \textbf{Answer:} Genetic studies
\\ \textit{Note:} Expected answer: Genetic studies
\item \textbf{Answer:} 40
\\ \textit{Note:} Expected answer: 40
\end{enumerate}

\vfill
\noindent\textit{End of Answer Key}
\end{document}